%==============================================================================
% 十、结论
%==============================================================================
\section{结论}

本文围绕商用啤酒风味代谢数据，依次完成感官评分预测、关键化合物归因与受约束的风味指纹重构三项任务，三问共用同一套数据核验与评估标准。

问题一以 231 维化学谱预测 50 个感官项，用外层 5 折、内层 5 折的嵌套交叉验证在岭回归、弹性网、LightGBM 与 ExtraTrees 四个模型间逐项选型。$R^2$ 中位 0.1469、均值 0.2223，49 项为正，可预测性最高的是酸度 0.7149 与酒精感 0.7108。0--6 分级由分位映射给出，加权 kappa 中位 0.3400、34 项超过 0.3；四组等级比例下 kappa 极差为 0.041，而“平移后直接取整”会使 53.2\% 的样本压到 0 分、kappa 跌至 0.1244。结论是分级映射不是收尾步骤，而是问题一能否成立的前提。

问题二先诊断附件2 的化合物网络：原谱聚类因为随机初始化退化为单模块（占 79\% 的节点），改用归一化拉普拉斯后模块度仍接近 0，而 Louvain 算得 $Q=0.3229$、高于度保持零分布的 95\% 分位，说明社区结构真实存在但与化学来源几乎无关（ARI 仅 0.0838），故归因以逐列置换重要性完成，化学族标签只用于解释。200 次自助重采样的入选频率用作稳定性度量，阈值由稳定选择的界反解为 0.5022，属操作性门槛。150 条 Top-3 记录中 19 条 OAV$>1$，其中同时越过筛选阈值的 10 条构成最强证据；另有 8 条接近阈值、40 条浓度远低于阈值、83 条因阈值缺项无法判断。跨模型 Jaccard 中位仅 0.0974，是本文必须如实披露的短板。

问题三以收缩高斯描述 22 个风格的可行域，最近邻距离比 1.076 说明生成没有退化为复制；用 8 成分偏最小二乘作 50 维感官代理，按问题二的 Top-3 并集与三个缺陷物把决策维压到 75，以 DMS、4VG 为硬约束求解 $\varepsilon$-约束的混合整数线性规划，300 个网格点全部可行、去重后得 7 个 Pareto 点。代理目标由基线 48.9620 降至 20.2063（降 58.7\%），前沿从改动 1 个化合物到 10 个连续可调；与手写 NSGA-II 的 115 个非支配点相比，精确前沿支配其中 92.17\%、反向支配为 0。低 $f_1$ 端的解依赖若干降到检出限的极端改动，若按每列 $\pm1$ 个 $\log_{10}$ 单位收紧可行域，降幅会明显收窄。跨风格迁移把偏差从起点 47.1450 降到 13.1366（降 72.14\%），改动了 25 个化合物，五类缺陷物全部不恶化。

三问之间不是并列关系。问题一给出“哪些感官项可预测、哪些不可”的分层判断，这个判断决定了问题二只能交付 19 条双条件交集而不是完整的 50×3 清单；问题二的 Top-3 并集又直接构成问题三的 75 维决策空间，也决定了问题三里哪些感官项能调、哪些只能被动接受。把这条链拆开各做一套，会得到三份互不相容的结论。

可以推广的三条经验：数据先拿官方对照表核一遍，不能照题面文字做；小样本下嵌套交叉验证不能省，否则报出的精度是挑出来的；结论要标可信范围，宁可否定自己的一部分结果，也不要给经不起追问的单一答案。
